% !TEX root = main.tex
\section{Limitations}
\label{sec:limitations}

The primary skill is exact three-digit addition under two-hidden-layer feedforward networks, extended in \texttt{experiments/archive/10\_autoregressive\_transformer\_replication/} and \texttt{experiments/archive/11\_transformer\_calibrated\_pretraining\_horizon/} to causal autoregressive transformers.
Nothing in the eleven studies estimates broad language-model benchmark difficulty or the fraction of benchmark variance that data would explain.
\texttt{experiments/archive/02\_heldout\_exposure\_difficulty\_laws/} states that limit in its own report, and the companion mathematics grid does not remove it.
The grid has 24 observations, a breadth proxy, and no measured starting difficulty.

The interventions identify specified training distributions.
Where digit marginals and answer histograms are unmatched, carrying is not an isolated cause.
Where they are matched, a residual input-digit total variation remains (mean 0.0398 in the carry-rate study).
We did not establish that the networks use the same internal procedure.
A fixed checker makes the scored skill stable.
It does not make the learned representation stable.

Several prespecified gates fail, and those failures constrain the claim.
Retention fails in the first study.
The starting-loss feature fails the 10\% line under shared warmups, and the unconstrained probe law in the answer-matched study diverges.
The constrained probe in the heterogeneous-history study loses to the unconstrained difficulty law by 0.82\%.
Online-code remediation misses the 85\% accuracy line because width 32 remains at a 27.80\% eight-seed mean.
The additive utility law worsens pooled error even though the same study's loss law succeeds.
\texttt{experiments/archive/09\_headroom\_bounded\_rlvr\_utility\_law/} caps that forecast by remaining headroom.
The cap is supported, at a relative reduction of 73.72\% [70.43, 77.26].
Probe features inside the cap are not established, at 8.10\% [2.52, 14.39], below the protocol's 10\% line.
A group-size comparison in the same folder is not established, at 0.96\% [0.44, 1.54].
Supported replications in that folder leave those two gates not established.
Near-ceiling widths shrink the accuracy reversal to a fraction of a point in one answer-matched replication.
A supported pooled gate can coexist with a small effect at large width.
\texttt{experiments/archive/10\_autoregressive\_transformer\_replication/} evaluates causal autoregressive transformers (2 layers, 4 attention heads, generating sum digits token-by-token).
No pooled gate is established there: the pure compositional reversal is 2.84 [2.01, 3.86] and 0.87 [0.62, 1.17] points against a 10-point line.
Eight-seed means at width 32 show a reversal, 11.76 points after mixture 0.1 but only 2.38 points after mixture 0.9, and remediation (93.19\% recovery; $+0.23$ points retention); these are not gate decisions.
However, transformers exhibit far greater sample-efficiency on arithmetic than feedforward MLPs: after 1,500 pretraining steps, widths 64, 96, 128, and 384 completely saturate at ceiling accuracy (>99.8\% to 100.0%, loss $<0.005$ bits) across both carry strata, leaving virtually zero residual error to predict.
Pooled held-out loss reduction (0.05\% [0.03, 2.84]) and prequential MDL reduction (3.77\% [1.21, 6.93]) remain below the 10\% gate, leaving all five prospective gates not established under the frozen criteria.
Transformers shift the saturation boundary, demonstrating that difficulty scaling laws require earlier checkpointing horizons or longer digit sequences to measure sub-ceiling dynamics.
\texttt{experiments/archive/11\_transformer\_calibrated\_pretraining\_horizon/} directly addresses that saturation boundary by calibrating the pretraining horizon to $T_{\text{pre}}=300$ updates.
Active headroom is restored across all five model widths ($d_{\text{model}} \in \{32, 64, 96, 128, 384\}$).
Pure compositional difficulty reversal is supported across all five widths (15.67 [12.13, 19.10] points at $q_{\text{pre}}=0.10$ vs 15.70 [12.89, 18.17] points at $q_{\text{pre}}=0.90$).
Prequential online MDL prediction is supported: probed starting difficulty cuts held-out online adaptation MDL root mean squared error by 18.84\% [5.88, 34.16].
The online-code reversal itself is not established: the pure two-carry cell costs 0.0666 [0.0606, 0.0715] extra bits per item after mixture 0.1, but the reverse gap after mixture 0.9 is $-0.0006$ [$-0.0075$, $0.0056$].
The online-code reversal therefore holds for the feedforward networks of Section~\ref{sec:mdl} but not for these transformers.
However, final conditional loss prediction is not established (0.99\% [$-11.40$, $5.66$]), because 1{,}000 updates of balanced post-training compress final test loss to 0.02 bits across all widths, whereas the online prequential code integrates over the active trajectory.
Additionally, while balanced remediation raises cascading two-carry accuracy from 68.32\% to 99.90\% (+31.58 points gain) and improves control accuracy by $+15.94$ [$+14.18$, $+17.80$] points, which passes the no-harm line, the remediation gate is not established because the lower bound of the initial gap (12.13 points) falls short of the protocol's 15-point minimum gap.

Conditional bits are not Kolmogorov complexity and, outside the online-code study, not a full minimum description length.
The online code charges for labels along one adaptation stream.
It does not charge for discovering the learning rule.
No study computes a shortest program.

Reproducibility has a numerical caveat that is separate from the scientific gates.
Training receipts for the later studies before \texttt{experiments/archive/09\_headroom\_bounded\_rlvr\_utility\_law/} record PyTorch 2.5.1+cu124 and a maximum replay difference of 0.0 bits.
A replay under PyTorch 2.7.0, reported on 10-17-2026 for the studies then in this draft, matched greedy predictions and produced bit differences on the order of $10^{-5}$.
The frozen assertion at $10^{-6}$ therefore fails on that build.
The first study's own receipt is PyTorch 2.7.0+cu126, with a maximum replay difference of 0.0 on its verification host.
We treat the $10^{-5}$ figure as a portability limitation, not as a change to any gate.
A 10-18-2026 replay of \texttt{experiments/archive/09\_headroom\_bounded\_rlvr\_utility\_law/} reproduced the written point estimates from the stored forecast table.
Greedy actions had 0 prediction mismatches on 800 final checkpoints, covering test and calibration, and on 200 warm probe checkpoints.
\texttt{test\_laws.py} passed.
Under PyTorch 2.7.0, \texttt{verify.py} stopped at a maximum bit difference of $7.62939453125\times 10^{-5}$.
That difference is above the frozen $10^{-6}$ line.
The training writeup for this study records PyTorch 2.5.1+cpu.
The frozen $10^{-6}$ assertion fails on that replay.
The tolerance stays at $10^{-6}$.

Coauthors and affiliation are unresolved.
Related-work coverage of benchmark contamination is intentionally absent, because this draft does not make that claim.
The verifiable-reward citation names the public phrase only.
\todocite{any additional related-work papers Brando wants cited after checking them against the claims above}
